\documentclass[12pt]{article}
\usepackage[utf8]{inputenc}
\usepackage{amsmath}
\usepackage{amssymb}

\title{Quantum Theories: The Missing Big Picture}
\author{A Unified Framework}
\date{}

\begin{document}

\maketitle

\section*{Blurb}

Quantum theory books abound, but readers are left in a fragmented maze. This course provides a unified framework by emphasizing quantum theories' goal: predict time evolution of states and observables.

\section*{The 11-Dimensional Framework}

Time evolution equations map onto 11-dimensional space:

1. Formalism (Operator vs Path Integral)
2. Relativity (Non-relativistic, Special, General)
3. Particle Number (Single, Fixed-N, Fock)
4. Picture (Schrodinger, Heisenberg, Interaction)
5. State Purity (Pure, Mixed)
6. System (Closed, Open)
7. Spin (Spinless, 1/2, Arbitrary)
8. Mass (Massive, Massless)
9. Approximation (Exact, Perturbative, Linearized)
10. Classical Limit (Quantum, Semiclassical, Classical)
11. Solutions (Analytical, Approximate, Numerical)

\section{Foundations: The Schrodinger Equation}

\subsection{Why Time Evolution?}

The central question of quantum mechanics is: How does a quantum system change with time?

Core Principle: All quantum theories answer through time evolution equations.

\[
i\hbar\frac{d}{dt}|\psi(t)\rangle = \hat{H}|\psi(t)\rangle
\]

\subsection{The 11 Dimensions}

D1: Formalism - How we express time evolution
D2: Relativity - Spacetime symmetries
D3: Particle Number - How many particles
D4: Picture - Where time dependence lives
D5: State Purity - Completeness of knowledge
D6: System - Boundary conditions
D7: Spin - Internal angular momentum
D8: Mass - Particle properties
D9: Approximation - Solution strategy
D10: Classical Limit - Quantum-classical transition
D11: Solutions - Implementation method

\section{The Simplest Case: Non-Relativistic Spinless}

\subsection{Elementary Schrodinger Equation}

\[
i\hbar\frac{\partial\psi(\mathbf{r},t)}{\partial t} = -\frac{\hbar^2}{2m}\nabla^2\psi(\mathbf{r},t) + V(\mathbf{r},t)\psi(\mathbf{r},t)
\]

Dimensions: Non-relativistic, Single particle, Schrodinger picture, Pure state, Closed system, Spinless, Massive.

\subsection{Operator Formalism}

Abstract form:
\[
i\hbar\frac{d|\psi(t)\rangle}{dt} = \hat{H}|\psi(t)\rangle
\]

where Hamiltonian:
\[
\hat{H} = \frac{\hat{p}^2}{2m} + V(\hat{\mathbf{r}},t)
\]

\subsection{Formal Solution}

For time-independent Hamiltonians:
\[
|\psi(t)\rangle = e^{-i\hat{H}t/\hbar}|\psi(0)\rangle
\]

The time evolution operator $U(t) = e^{-i\hat{H}t/\hbar}$ is unitary and preserves probability.

\subsection{Energy Eigenstates}

When V is time-independent, expand in energy eigenstates:
\[
\hat{H}|n\rangle = E_n|n\rangle
\]

yielding:
\[
|\psi(t)\rangle = \sum_n c_n e^{-iE_nt/\hbar}|n\rangle
\]

where $c_n = \langle n|\psi(0)\rangle$.

\subsection{The Hydrogen Atom}

Hamiltonian for electron orbiting proton:
\[
\hat{H} = \frac{\hat{p}^2}{2m_e} - \frac{ke^2}{r}
\]

Exact solution yields energy levels:
\[
E_n = -\frac{13.6 \text{ eV}}{n^2}, \quad n = 1,2,3,\ldots
\]

\section{Mathematical Foundations}

\subsection{Hilbert Space}

Quantum states live in complex Hilbert space with:

- Inner product: $\langle\phi|\psi\rangle \in \mathbb{C}$
- Linearity of operators
- Completeness of basis

\subsection{Hermitian Operators}

Observables: $\hat{A}^\dagger = \hat{A}$

Properties:
- Real eigenvalues
- Orthogonal eigenstates
- Complete basis

\subsection{Canonical Commutation Relations}

\[
[\hat{x}_i, \hat{p}_j] = i\hbar\delta_{ij}
\]

Heisenberg uncertainty principle:
\[
\Delta x \cdot \Delta p \geq \frac{\hbar}{2}
\]

\section{Dimension Variations: Pictures}

\subsection{Heisenberg Picture}

States static, operators evolve:
\[
|\psi_H(t)\rangle = |\psi_H(0)\rangle
\]

Operators evolve as:
\[
\hat{A}_H(t) = \hat{U}^\dagger(t)\hat{A}_S\hat{U}(t)
\]

where $\hat{U}(t) = e^{-i\hat{H}t/\hbar}$.

\subsection{Interaction Picture}

For perturbations where $\hat{H} = \hat{H}_0 + \hat{H}_I(t)$:

States and operators both evolve relative to unperturbed Hamiltonian.

\section{Relativistic Quantum Mechanics}

\subsection{Klein-Gordon Equation}

From $E^2 = p^2c^2 + m^2c^4$:

\[
\left(\frac{1}{c^2}\frac{\partial^2}{\partial t^2} - \nabla^2 + \frac{m^2c^2}{\hbar^2}\right)\phi(\mathbf{r},t) = 0
\]

Describes spin-0 particles.

\subsection{Dirac Equation}

Incorporating spin-1/2:

\[
i\hbar\frac{\partial\psi}{\partial t} = (\boldsymbol{\alpha} \cdot \mathbf{p}c + \beta mc^2)\psi
\]

Solutions naturally incorporate antiparticles.

\section{Quantum Field Theory}

\subsection{From Particles to Fields}

Second quantization leads to quantum field theory.

Field operators:
\[
\hat{\phi}(\mathbf{r}) = \int \frac{d^3p}{(2\pi)^{3/2}\sqrt{2E_\mathbf{p}}}(\hat{a}_\mathbf{p}e^{i\mathbf{p}\cdot\mathbf{r}} + \hat{a}_\mathbf{p}^\dagger e^{-i\mathbf{p}\cdot\mathbf{r}})
\]

Creation and annihilation operators satisfy:

$[\hat{a}_p, \hat{a}_p^\dagger] = 1$ (bosons)

$\{\hat{c}_p, \hat{c}_p^\dagger\} = 1$ (fermions)

\section{Quantum Electrodynamics}

\subsection{Electron-Photon Coupling}

Lagrangian:
\[
\mathcal{L} = \bar{\psi}(i\gamma^\mu(\partial_\mu - ieA_\mu) - m)\psi - \frac{1}{4}F_{\mu\nu}F^{\mu\nu}
\]

Coupling constant: $\alpha \approx 1/137$

QED predictions:
- Lamb shift: 1.06 GHz (theory) vs 1.057 GHz (experiment)
- Electron g-factor: agreement to 1 part in $10^{12}$

\section{Open Systems}

\subsection{Density Matrix}

For mixed states:
\[
\rho = \sum_i p_i |\psi_i\rangle\langle\psi_i|
\]

Expectation value:
\[
\langle A\rangle = \text{Tr}[\rho\hat{A}]
\]

\subsection{Master Equations}

Lindblad master equation:
\[
\frac{d\rho_S}{dt} = -\frac{i}{\hbar}[\hat{H}_S, \rho_S] + \sum_\mu\left(\hat{L}_\mu\rho_S\hat{L}_\mu^\dagger - \frac{1}{2}\{\hat{L}_\mu^\dagger\hat{L}_\mu, \rho_S\}\right)
\]

\section{Solution Techniques}

\subsection{Perturbation Theory}

Time-independent with $\hat{H} = \hat{H}_0 + \lambda\hat{H}_1$:

\[
E_n = E_n^{(0)} + \lambda E_n^{(1)} + \lambda^2 E_n^{(2)} + \ldots
\]

First-order correction:
\[
E_n^{(1)} = \langle n^{(0)}|\hat{H}_1|n^{(0)}\rangle
\]

\subsection{Path Integrals}

Feynman path integral:
\[
\langle x_f, t_f|x_i, t_i\rangle = \int \mathcal{D}x \, e^{iS[x]/\hbar}
\]

where action:
\[
S[x] = \int_{t_i}^{t_f} dt \left(\frac{m}{2}\dot{x}^2 - V(x)\right)
\]

\subsection{WKB Approximation}

For slowly varying potential:
\[
\psi(x) = \frac{A}{\sqrt{p(x)}}e^{\pm i\int p(x)dx/\hbar}
\]

Tunneling probability:
\[
T \approx e^{-2\gamma}
\]

\subsection{Numerical Methods}

Imaginary time evolution to find ground state:
\[
\frac{\partial\psi}{\partial\tau} = -\hat{H}\psi
\]

where $\tau = it$.

\section{Connecting the Dimensions}

\subsection{Taxonomy of Quantum Theories}

Quantum Mechanics: Non-rel, Single, Pure - Atoms, molecules

QFT: Relativistic, Fock - Particle physics

QED: Special rel, Massive - Electromagnetic

QCD: Special rel, Massless - Strong force

Open Systems: Open, Mixed - Decoherence

\subsection{The Schrodinger Equation as Lodestar}

Universal across all theories:
\[
i\hbar\frac{d}{dt}|\psi(t)\rangle = \hat{H}(t)|\psi(t)\rangle
\]

Variations:
- Different Hamiltonians
- Different Hilbert spaces
- Different representations
- Different solution methods

\section{Open Problems and Conclusion}

\subsection{Quantum Gravity}

Quantize gravity without pathological infinities.

Approaches: String Theory, Loop Quantum Gravity, Causal Sets

\subsection{From Fragmentation to Unity}

Quantum mechanics appears as disconnected topics: Schrodinger equation, Dirac equation, Path integrals, Quantum field theory, Quantum computing.

The unifying principle: time evolution of states and observables.

The 11-dimensional framework provides:

1. Conceptual scaffolding
2. Predictive power
3. Problem-solving tools
4. Research guidance

\subsection{Final Thought}

Quantum mechanics is not mysterious. The missing big picture is that quantum theories form a coherent whole once viewed through time evolution in 11-dimensional space.

The Schrodinger equation is the starting point from which all quantum theories emerge through systematic variation of physical assumptions.

\textit{The quantum world awaits those with the right framework to see it clearly.}

\section*{Appendix: Key Equations}

Schrodinger: $i\hbar\frac{\partial\psi}{\partial t} = \hat{H}\psi$

Commutation: $[\hat{x}_i, \hat{p}_j] = i\hbar\delta_{ij}$

Uncertainty: $\Delta x \cdot \Delta p \geq \hbar/2$

Klein-Gordon: $(\Box + m^2)\phi = 0$

Dirac: $(i\gamma^\mu\partial_\mu - m)\psi = 0$

Partition Function: $Z = \text{Tr}[e^{-\beta\hat{H}}]$

\end{document}
